% MAS1004 Assignment 2: long report.
%
% You fill this in with your agent, in your repository. AGENTS.md tells the
% agent how. Red text in the PDF is what still has to be replaced.
%
% - Pictures go in figures/. Copy them from results/, which is not in git.
% - Outputs go in outputs/: copy experiments.csv, LOG.md and the outputs of
%   check_submission.py and curate.py count here. \outputfile{outputs/...}
%   prints a file exactly as it is. Every number in this report has to come
%   from such an output.
%
% To make the PDF: compress this folder into a zip, and on Overleaf choose
% New Project, Upload Project. Set Menu, Compiler to XeLaTeX, and Recompile.

\documentclass[11pt]{article}
\usepackage{mas1004}

\reporttitle{Assignment 2: Long report}
\studentname{\fillin{your name}}
\studentid{\fillin{your student ID}}

\begin{document}
\makeheader

Leaderboard nickname: \fillin{your nickname}. My last release of each track is
\texttt{track-a-1} (Track A: the starter recipe on all training photos, from
commit \texttt{ff34d66}) and \texttt{track-b-1} (Track B: the model called b1).
The leaderboard scores next to my validation scores are a note for tomorrow:
\fillin{the leaderboard score of \texttt{track-a-1} and of \texttt{track-b-1}
next to the validation scores 0.4942 / 0.4971 (Track A) and 0.4835 (Track B),
read from the leaderboard in the morning --- not written yet}.

% The sections below are a suggestion. Change, merge or add sections as you
% like; README.md, "The long report", says what a reader should be able to see.

% ---------------------------------------------------------------------------
\section*{Track A}

\begin{center}\footnotesize
\begin{tabular}{lp{0.55\linewidth}r}
\toprule
tag & what changed from the run before & validation \\
\midrule
\texttt{baseline}, seed 0 & nothing: the starter recipe on all 44,257 training photos & 0.4942 \\
\texttt{baseline}, seed 1 & the same recipe and the same photos; only the seed changed & 0.4971 \\
\texttt{no\_flip}, seed 0 & the left-right flip removed from \texttt{train\_transform} & 0.4926 \\
\texttt{drop\_nonstreet}, seed 0 & the flip back, so the recipe is the starter recipe again, on a selection of 44,172 photos instead of 44,257: 85 dropped as not a street scene & 0.4891 \\
\bottomrule
\end{tabular}
\end{center}

Every number in that table is from \texttt{experiments.csv}, which is printed in
full under ``All runs''. All four runs went through \texttt{src/official.py} with
\texttt{rules.TRACK\_A\_BUDGET}, 249,984 photos each, on Colab's GPU.

\textbf{Why I made each change, and what I think the result means} (my own words,
from \texttt{LOG.md}):

\textbf{baseline, seeds 0 and 1.} Why: to have a reference, and to see how much
the score moves by chance when only the seed changes. What it means: the two
seeds differ by about 0.003, so a change has to move the score by more than that
before I believe it.

\textbf{no\_flip.} Why: a flip turns signs into mirror writing and puts traffic
on the wrong side, and Japan drives on the left while Korea and Taiwan drive on
the right, so I thought it destroys clues. What it means: 0.4926 is within the
seed noise of baseline, so removing the flip did not help; maybe signs and
traffic direction are a small part of most photos.

\textbf{drop\_nonstreet.} Why: my rule, written before I chose which photos to
look at --- drop training photos you could not see while walking a street:
nature only, inside a room (museum, restaurant, office), exhibits, food, maps, a
plane in the sky; keep walkable covered spaces such as shopping arcades,
underground malls and stations; street photos that are just hard to tell from
another city stay in. What it means: as I predicted, removing 85 photos
(0.2\%) did not help. It came out about 0.005 below baseline, but with only one
seed and per-city swings of 0.02--0.04 between seeds, I cannot tell whether that
is a real drop or noise. To see a real effect I would have to apply the rule to
all training photos, not only the 48 most doubtful per city.

\textbf{My final recipe and selection.} The recipe is the starter recipe: in
\texttt{experiments.csv} the recipe hash of \texttt{drop\_nonstreet}
(\texttt{17ac485223fa}) is the same as both \texttt{baseline} runs, because the
left-right flip is back and nothing else in \texttt{track\_a/recipe.py} differs.
The selection is \texttt{track\_a/selection.csv}, 44,172 of the 44,257 training
photos, which validated at 0.4891. The run that is actually released is
\texttt{track-a-1}, the starter recipe on all training photos, with validation
0.4942 (seed 0) and 0.4971 (seed 1). Its leaderboard score next to those two
numbers: \fillin{to be filled in from tomorrow's leaderboard --- not written
yet}.

% ---------------------------------------------------------------------------
\section*{Decisions you made from photos}

Two kinds of sheets were drawn from the \texttt{baseline} seed 0 model on Colab,
both numbered so that a photo can be named by its sheet and its number:
\texttt{curate.py mistakes} (\texttt{-{}-every 2}, 6,039 validation photos scored,
showing the ones the model got wrong most confidently) and
\texttt{curate.py suspects} (\texttt{-{}-every 4}, 11,065 training photos scored,
showing the 48 photos per city with the lowest probability for their own city).
All the counts in this section are from those two runs' output in
\texttt{results/colab\_logs/}.

\picturefile[0.82\linewidth]{figures/mistakes_baseline_seed0_busan_00.jpg}

\begin{center}\footnotesize
Busan, \texttt{mistakes} sheet: the 48 validation photos baseline seed 0 got
wrong most confidently. Of the 695 Busan validation photos read with
\texttt{-{}-every 2}, 430 were wrong.
\end{center}

\sidebyside{figures/suspects_baseline_seed0_busan_00.jpg}{Busan,
\texttt{suspects} sheet: the 48 Busan training photos with the lowest Busan
probability. I dropped 24 of the 48 on this sheet.}
{figures/suspects_baseline_seed0_tokyo_00.jpg}{Tokyo, \texttt{suspects} sheet:
the same for Tokyo. I dropped 22 of the 48 on this sheet.}

\textbf{The rule I wrote before choosing which photos to drop.} Drop training
photos you could not see while walking a street --- nature only, inside a room
(museum, restaurant, office), exhibits, food, maps, a plane in the sky. Keep
walkable covered spaces such as shopping arcades, underground malls and
stations. Street photos that are just hard to tell from another city stay in.

I applied that rule to the five \texttt{suspects} sheets of \texttt{baseline}
seed 0: osaka, tokyo, seoul, busan and taipei.

\textbf{What it changed.} 85 photos left the selection (osaka 14, tokyo 22,
seoul 12, busan 24, taipei 13, which is the difference between the two columns
of the count below), so \texttt{track\_a/selection.csv} went from
44,257 rows to 44,172 rows. Every dropped photo and the reason are in
\texttt{track\_a/dropped.csv}. The photos are not deleted, only left out, so
\texttt{curate.py undrop} can put any of them back. The scores before and after:
0.4942 (baseline, seed 0) and 0.4971 (baseline, seed 1) before, 0.4891 after,
with the same recipe. The output of \texttt{curate.py count} on the new
selection:

\outputfile{outputs/curate_count.txt}

\textbf{The requests I gave my agent for these decisions} (step 7 of my work
with it, copied exactly as I sent it):

\begin{Verbatim}[breaklines,breakanywhere,fontsize=\footnotesize,frame=single]
My rule, written before choosing: drop training photos you could not see while walking a street — nature only, inside a room (museum, restaurant, office), exhibits, food, maps, a plane in the sky. Keep walkable covered spaces such as shopping arcades, underground malls and stations. Street photos that are just hard to tell from another city stay in. I applied it to all five suspect sheets of baseline seed 0.

Drop photos 1, 6, 10, 11, 20, 33, 34, 36, 37, 39, 41, 46, 47, 48 from suspects_baseline_seed0_osaka_00, reason "not a street scene".
Drop photos 1, 6, 8, 9, 12, 13, 14, 16, 17, 18, 26, 28, 30, 31, 32, 34, 38, 39, 42, 43, 45, 46 from suspects_baseline_seed0_tokyo_00, reason "not a street scene".
Drop photos 4, 12, 14, 20, 21, 30, 36, 37, 39, 41, 43, 45 from suspects_baseline_seed0_seoul_00, reason "not a street scene".
Drop photos 1, 2, 6, 7, 8, 10, 12, 17, 21, 24, 26, 27, 29, 30, 33, 35, 37, 38, 39, 44, 45, 46, 47, 48 from suspects_baseline_seed0_busan_00, reason "not a street scene".
Drop photos 1, 2, 4, 10, 11, 14, 15, 19, 21, 32, 38, 47, 48 from suspects_baseline_seed0_taipei_00, reason "not a street scene".

Then put the left-right flip back in track_a/recipe.py so the recipe is the same as baseline, change nothing else, and show me the diff. Train the new selection on Colab with the tag drop_nonstreet, and tell me when to press Run all.
\end{Verbatim}

% ---------------------------------------------------------------------------
\section*{Track B}

\textbf{b1, the first step.} \texttt{train\_b.py} as it came: \texttt{resnet18}
with its ImageNet weights, 224 px, 4 epochs, batch 64, no balancing, on every
photo in \texttt{data/starter/train}, submitted with \texttt{-{}-submit}. It
validated at 0.4835 (seoul 0.591, busan 0.3791, tokyo 0.5986, osaka 0.3369,
taipei 0.5118) after 176,896 photos through the network, from the b1 row of
\texttt{experiments.csv}; the run's own output is in
\texttt{results/colab\_logs/}.

Why: to have a valid Track B submission and a reference for later Track B
changes. What it means: b1 uses the same model, image size and augmentation as
the Track A baseline, but 4 epochs is only about 177,000 photos instead of
250,000. Its score is about 0.01 below baseline, more than the 0.003 seed noise,
so I think it is lower mainly because it trained on fewer photos.

\textbf{b2 did not run.} The next step was the same file with 6 epochs instead
of 4 --- \texttt{python src/train\_b.py -{}-tag b2 -{}-epochs 6 -{}-submit} followed
by \texttt{python src/check\_submission.py -{}-no-val}, queued in
\texttt{colab\_runs.txt}. The run never started: the free Colab GPU limit was
reached, so Colab would not give me a GPU that day. There is therefore no b2 row
in \texttt{experiments.csv}, no b2 validation score, and \texttt{submission/}
still holds b1. Nothing else about b2 was changed: \texttt{train\_b.py} is
identical to the file b1 used, and 6 epochs instead of 4 is the only difference.

\textbf{Extra photos.} None. Both Track B runs used only
\texttt{data/starter/train}, so there was nothing to collect from a phone or put
in Google Drive.

\textbf{My validation score next to my leaderboard score.} b1 validated at
0.4835, and \texttt{track-b-1} is the released model.
\fillin{its leaderboard score next to 0.4835, and why the two might disagree, in
my own words --- to be read from the leaderboard tomorrow and written then}.

\subsection*{Your submission, as the grader sees it}
\outputfile{outputs/check_submission.txt}

The description of the model that \texttt{train\_b.py} wrote next to it, in
\texttt{submission/submission.json}:

\outputfile{outputs/submission.json}

% ---------------------------------------------------------------------------
\section*{What made the difference}

Neither Track A change I measured beat the baseline. With the starter recipe on
all training photos the score was 0.4942 (seed 0) and 0.4971 (seed 1); removing
the left-right flip gave 0.4926, and taking the 85 non-street photos out of the
selection gave 0.4891. Track B never got past its reference run: b1 validated at
0.4835 on 176,896 photos, and b2, the same file with 6 epochs, could not run
because the free Colab GPU limit was reached. Every number here is a
\texttt{validation} value from \texttt{experiments.csv}.

\fillin{Which change helped most and which least, in my own words, and what I
would do next with one more week --- not written yet.}

My plan for Track A, from \texttt{LOG.md}, is to apply the photo rule to all
training photos, not only the 48 most doubtful per city, because 85 photos out
of 44,257 was too small a change to see through the seed noise. The 6-epoch
Track B run for b2 is still queued in \texttt{colab\_runs.txt} and has not run
yet.

% ---------------------------------------------------------------------------
\section*{All runs}
\outputfile{outputs/experiments.csv}

\section*{Lab notebook}
\outputfile{outputs/LOG.md}

\end{document}
